\documentclass[11pt,a4paper]{article}

\usepackage[utf8]{inputenc}
\usepackage[T1]{fontenc}
\usepackage[italian]{babel}
\usepackage[margin=25mm]{geometry}
\usepackage{amsmath,amssymb}
\usepackage{siunitx}
\usepackage{booktabs}
\usepackage{listings}
\usepackage{xcolor}
\usepackage{tikz}
\usetikzlibrary{positioning,arrows.meta,fit,backgrounds,calc,shapes.geometric}
\usepackage{hyperref}

\sisetup{detect-all, output-decimal-marker={,}, per-mode=symbol}

\lstset{
  basicstyle=\ttfamily\small,
  frame=single,
  breaklines=true,
  columns=fullflexible,
  keepspaces=true
}

\title{Stima dell'invecchiamento di un azionamento a IGBT\\
\large Modello di Arrhenius con soglie congiunte di corrente e temperatura acqua}
\author{}
\date{\today}

\begin{document}
\maketitle

\section{Scopo e ipotesi}

Il documento definisce un modello empirico per la stima dell'invecchiamento
cumulato di un azionamento raffreddato a liquido, a partire da grandezze
disponibili in linea:

\begin{itemize}
  \item corrente erogata $I$ (valore efficace di fase);
  \item temperatura dell'acqua in ingresso al dissipatore $T_w$;
  \item temperatura di giunzione del modulo IGBT $T_j$ (misurata o stimata);
  \item numero di interventi di manutenzione/riparazione $N$ effettuati
        sull'azionamento;
  \item tipo di modulo IGBT, $\tau \in \{1,2\}$.
\end{itemize}

L'impostazione segue l'approccio \emph{physics-of-failure} adottato in
letteratura per l'affidabilità dei convertitori di
potenza~\cite{wang2013,ma2016,lutz2018}.

Ipotesi di base:

\begin{enumerate}
  \item il meccanismo di degrado dominante è di tipo termicamente attivato,
        quindi descrivibile con una legge di Arrhenius sulla temperatura di
        giunzione;
  \item la corrente agisce come fattore di stress moltiplicativo (perdite,
        densità di corrente nei bonding wire, ampiezza dei cicli termici);
  \item esiste una condizione di funzionamento \emph{fuori progetto} che si
        verifica quando la temperatura dell'acqua supera una soglia
        \emph{contemporaneamente} al superamento di una soglia di corrente:
        in tale condizione il tasso di invecchiamento è ulteriormente
        accelerato;
  \item il modulo di tipo~1 ha coefficiente di invecchiamento doppio rispetto
        al tipo~2;
  \item ogni intervento di manutenzione lascia una penalizzazione residua
        (rimontaggio, stress termomeccanico sul pacco, degrado dell'interfaccia
        termica).
\end{enumerate}

\section{Notazione}

\begin{table}[h]
\centering
\begin{tabular}{llc}
\toprule
Simbolo & Significato & Unità \\
\midrule
$I$            & corrente erogata (RMS)                              & \si{\ampere} \\
$I_N$          & corrente nominale dell'azionamento                  & \si{\ampere} \\
$I_{th}$       & soglia di corrente per il regime accelerato         & \si{\ampere} \\
$T_w$          & temperatura acqua in ingresso                       & \si{\degreeCelsius} \\
$T_{w,th}$     & soglia di temperatura acqua                         & \si{\degreeCelsius} \\
$T_j$          & temperatura di giunzione IGBT                       & \si{\kelvin} \\
$T_{ref}$      & temperatura di riferimento del modello              & \si{\kelvin} \\
$E_a$          & energia di attivazione                              & \si{\electronvolt} \\
$k_B$          & costante di Boltzmann, \num{8.617e-5}               & \si{\electronvolt\per\kelvin} \\
$m$            & esponente di stress di corrente                     & -- \\
$\beta$        & guadagno di accelerazione in regime di soglia       & -- \\
$\gamma$       & penalizzazione per intervento di manutenzione       & -- \\
$k_\tau$       & coefficiente di invecchiamento del tipo di IGBT     & -- \\
$L_0$          & vita attesa nelle condizioni di riferimento         & \si{\hour} \\
$A$            & indice di invecchiamento cumulato, $A\in[0,1]$      & -- \\
\bottomrule
\end{tabular}
\caption{Simboli utilizzati.}
\end{table}

\section{Schema complessivo del modello}
\label{sec:schema}

La Figura~\ref{fig:schema} riassume la struttura del modello: le grandezze
misurate a sinistra, i fattori di accelerazione al centro, l'accumulatore di
danno e le uscite diagnostiche a destra.

\begin{figure}[htbp]
\centering
\resizebox{\textwidth}{!}{%
\begin{tikzpicture}[
  font=\small,
  inp/.style={draw, rounded corners, fill=blue!8, minimum width=34mm,
              minimum height=9mm, align=center},
  blk/.style={draw, fill=orange!12, minimum width=40mm,
              minimum height=10mm, align=center},
  outb/.style={draw, rounded corners, fill=green!12, minimum width=40mm,
              minimum height=10mm, align=center},
  op/.style={draw, circle, fill=gray!18, inner sep=0.5pt, minimum size=8mm},
  ar/.style={-{Stealth[length=2.2mm]}, thick},
  dar/.style={-{Stealth[length=2.2mm]}, thick, dashed, gray!70!black}
]

%% ---- ingressi ----
\node[inp] (I)   at (0, 4.6) {$I$\\corrente erogata};
\node[inp] (Tw)  at (0, 3.2) {$T_w$\\temp. acqua ingresso};
\node[inp] (Tj)  at (0, 1.8) {$T_j$\\temp. giunzione IGBT};
\node[inp] (N)   at (0, 0.4) {$N$\\n.\ interventi};
\node[inp] (TY)  at (0,-1.0) {$\tau$\\tipo di modulo};

%% ---- fattori ----
\node[blk] (AFI)  at (6.0, 4.6) {Stress di corrente\\$AF_I=(I/I_N)^m$};
\node[blk] (LAM)  at (6.0, 3.2) {Soglia congiunta\\$1+\beta\,\Lambda_s(I,T_w)$};
\node[blk] (ARR)  at (6.0, 1.8) {Arrhenius\\$e^{\frac{E_a}{k_B}\left(\frac{1}{T_{ref}}-\frac{1}{T_j}\right)}$};
\node[blk] (AFN)  at (6.0, 0.4) {Manutenzione\\$1+\gamma N$};
\node[blk] (KT)   at (6.0,-1.0) {Tipo IGBT\\$k_1=2k_2$};

%% ---- prodotto e uscite ----
\node[op]  (MUL)  at (10.4, 1.8) {$\times$};
\node[blk, fill=red!10] (RATE) at (14.6, 1.8)
     {Tasso di danno\\$\dfrac{dA}{dt}=\dfrac{k_\tau}{L_0}\prod AF_i$};
\node[outb] (INT)  at (14.6, 0.0) {Accumulatore\\$A=\int \frac{dA}{dt}\,dt$};
\node[outb] (RES)  at (14.6,-1.8) {Vita residua\\$L_{res}=(1-A)\big/\frac{dA}{dt}$};

%% ---- collegamenti ----
\draw[ar] (I)  -- (AFI);
\draw[ar] (Tw) -- (LAM);
\draw[ar] (Tj) -- (ARR);
\draw[ar] (N)  -- (AFN);
\draw[ar] (TY) -- (KT);
\draw[ar] (I.east) -- ++(0.6,0) |- (LAM.west);
\draw[dar] ($(Tw.south)+(1.15,0)$) -- ($(Tj.north)+(1.15,0)$)
      node[midway, right=1pt, font=\scriptsize, align=left]
      {stima $T_j$\\eq.~\eqref{eq:termico}};

\foreach \n in {AFI,LAM,ARR,AFN,KT}{
  \draw[ar] (\n.east) -- ++(0.7,0) |- (MUL);
}
\draw[ar] (MUL) -- (RATE);
\draw[ar] (RATE) -- (INT);
\draw[ar] (INT) -- (RES);

%% ---- riquadri esplicativi ----
\node[align=center, font=\footnotesize\itshape] at (0,-2.4) {grandezze misurate\\o registrate};
\node[align=center, font=\footnotesize\itshape] at (6.0,-2.4) {fattori di accelerazione};
\node[align=center, font=\footnotesize\itshape] at (14.6,-3.2) {uscite diagnostiche};

\end{tikzpicture}}
\caption{Struttura del modello di stima dell'invecchiamento: gli ingressi
misurati generano fattori di accelerazione moltiplicativi che, combinati con il
coefficiente di tipo modulo $k_\tau$, definiscono il tasso di consumo di vita;
l'integrazione nel tempo fornisce l'indice cumulato e la vita residua.}
\label{fig:schema}
\end{figure}

La Figura~\ref{fig:soglie} mostra il piano operativo $(I,T_w)$ con il quadrante
critico in cui si attiva il fattore di accelerazione aggiuntivo: l'attivazione
richiede il superamento \emph{simultaneo} delle due soglie.

\begin{figure}[htbp]
\centering
\begin{tikzpicture}[font=\small, scale=1.0]

% area critica
\fill[red!15] (4.2,3.0) rectangle (8.0,5.4);
% aree non critiche
\fill[green!10] (0,0) rectangle (4.2,3.0);
\fill[yellow!15] (4.2,0) rectangle (8.0,3.0);
\fill[yellow!15] (0,3.0) rectangle (4.2,5.4);

% griglia soglie
\draw[dashed, thick] (4.2,0) -- (4.2,5.4);
\draw[dashed, thick] (0,3.0) -- (8.0,3.0);

% assi
\draw[-{Stealth[length=2.5mm]}, thick] (0,0) -- (8.6,0) node[below left=2pt] {$I$};
\draw[-{Stealth[length=2.5mm]}, thick] (0,0) -- (0,5.9) node[below left=2pt, rotate=90] {$T_w$};

\node[below] at (4.2,-0.1) {$I_{th}=0.8\,I_N$};
\node[left]  at (-0.1,3.0) {$T_{w,th}$};

\node[align=center, font=\small] at (2.1,1.5) {$\Lambda=0$\\normale};
\node[align=center, font=\small] at (6.1,1.5) {$\Lambda=0$};
\node[align=center, font=\small] at (2.1,4.2) {$\Lambda=0$};
\node[align=center, font=\small\bfseries] at (6.1,4.2)
     {$\Lambda=1$\\regime critico\\$AF_{th}=1+\beta$};

% transizione regolarizzata
\draw[<->, gray!70!black] (4.2,5.65) -- (4.7,5.65)
      node[midway, above, font=\scriptsize] {$\Delta I$};
\draw[<->, gray!70!black] (8.35,3.0) -- (8.35,3.5)
      node[midway, right, font=\scriptsize] {$\Delta T_w$};

\end{tikzpicture}
\caption{Piano operativo corrente--temperatura acqua. Il fattore di
accelerazione aggiuntivo si attiva solo nel quadrante in alto a destra;
$\Delta I$ e $\Delta T_w$ sono le ampiezze di transizione della
formulazione regolarizzata~\eqref{eq:lambda_soft}.}
\label{fig:soglie}
\end{figure}

\section{Modello termico semplificato}

Se $T_j$ non è disponibile da misura diretta (NTC nel modulo, $V_{ce,on}$ come
parametro elettrico termosensibile~\cite{avenas2012}), si può stimare in regime quasi-stazionario a partire da $T_w$ e dalle
perdite:

\begin{equation}
  T_j = T_w + R_{th,jw}\,P_{loss}(I),
  \qquad
  P_{loss}(I) \simeq a_1 I + a_2 I^2 ,
  \label{eq:termico}
\end{equation}

dove il termine lineare rappresenta le perdite di commutazione (approssimate
come proporzionali alla corrente commutata) e il termine quadratico le perdite
di conduzione sulla resistenza differenziale. $R_{th,jw}$ include giunzione--case,
interfaccia termica e case--acqua, ed è funzione della portata del circuito di
raffreddamento.

\section{Fattore di accelerazione di Arrhenius}

Il fattore di accelerazione rispetto alla temperatura di riferimento vale

\begin{equation}
  AF_{Arr}(T_j) = \exp\!\left[\frac{E_a}{k_B}
  \left(\frac{1}{T_{ref}} - \frac{1}{T_j}\right)\right].
  \label{eq:arrhenius}
\end{equation}

Per $T_j > T_{ref}$ si ha $AF_{Arr} > 1$ (invecchiamento accelerato), per
$T_j < T_{ref}$ si ha $AF_{Arr} < 1$. Le temperature vanno espresse in kelvin.
Valori tipici di $E_a$ per moduli di potenza si collocano fra
\SI{0.4}{\electronvolt} e \SI{0.8}{\electronvolt} a seconda del meccanismo
dominante (degrado metallizzazione, elettromigrazione, degrado del
gel)~\cite{ciappa2002,jep122}.

\section{Fattore di stress di corrente}

\begin{equation}
  AF_I(I) = \left(\frac{I}{I_N}\right)^{m},
  \qquad m \in [1.5,\,4].
  \label{eq:corrente}
\end{equation}

L'esponente $m$ raccoglie in forma compatta la dipendenza del danno dalle
escursioni termiche di ciclo (tipo Coffin--Manson, $N_f \propto \Delta T_j^{-n}$,
con $\Delta T_j$ funzione crescente di $I$); i modelli di riferimento per il
power cycling dei moduli IGBT sono quelli di
LESIT~\cite{held1997}, CIPS08~\cite{bayerer2008} e
Scheuermann~\cite{scheuermann2013}.

\section{Condizione di soglia congiunta}

Il regime di invecchiamento accelerato si attiva quando \emph{entrambe} le
condizioni sono verificate:

\begin{equation}
  \Lambda(I,T_w) =
  \mathbb{1}\!\left[T_w > T_{w,th}\right]\cdot
  \mathbb{1}\!\left[I > I_{th}\right] \in \{0,1\}.
  \label{eq:lambda_hard}
\end{equation}

La formulazione a gradino \eqref{eq:lambda_hard} è discontinua e, in presenza di
rumore di misura, produce commutazioni ripetute dell'indice. Si preferisce
quindi una versione regolarizzata con funzioni sigmoidali:

\begin{equation}
  \Lambda_s(I,T_w) =
  \sigma\!\left(\frac{T_w - T_{w,th}}{\Delta T_w}\right)\cdot
  \sigma\!\left(\frac{I - I_{th}}{\Delta I}\right),
  \qquad
  \sigma(x) = \frac{1}{1+e^{-x}} ,
  \label{eq:lambda_soft}
\end{equation}

con $\Delta T_w$ e $\Delta I$ ampiezze di transizione (ad esempio
\SI{2}{\kelvin} e $0.05\,I_N$). In alternativa si può mantenere la forma a
gradino aggiungendo isteresi sulle due soglie.

Il fattore di accelerazione complessivo della condizione fuori progetto è

\begin{equation}
  AF_{th}(I,T_w) = 1 + \beta\,\Lambda_s(I,T_w),
\end{equation}

con $\beta = 1$ che corrisponde al raddoppio del tasso di invecchiamento nel
regime critico.

\section{Tipo di IGBT}

\begin{equation}
  k_\tau =
  \begin{cases}
    2\,k_0 & \tau = 1 \quad \text{(modulo di tipo 1)}\\[2pt]
    k_0    & \tau = 2 \quad \text{(modulo di tipo 2)}
  \end{cases}
  \qquad k_0 = 1 .
  \label{eq:tipo}
\end{equation}

A parità di profilo di missione, il modulo di tipo~1 consuma vita al doppio
della velocità del tipo~2.

\section{Interventi di manutenzione}

Ogni intervento sull'azionamento introduce una penalizzazione permanente
modellata come fattore moltiplicativo lineare:

\begin{equation}
  AF_N(N) = 1 + \gamma N ,
  \qquad \gamma \in [0.02,\,0.10].
  \label{eq:interventi}
\end{equation}

Il conteggio $N$ è incrementato a ogni intervento registrato (sostituzione
moduli, apertura del circuito di raffreddamento, rimontaggio del pacco di
potenza). Se si vuole distinguere la gravità degli interventi, si sostituisce
$\gamma N$ con $\sum_{i=1}^{N}\gamma_i$, con $\gamma_i$ pesato per tipologia.

\section{Formula complessiva}

Il tasso istantaneo di consumo di vita è

\begin{equation}
  \boxed{\;
  \frac{dA}{dt} =
  \frac{k_\tau}{L_0}\,
  \left(\frac{I}{I_N}\right)^{m}
  \exp\!\left[\frac{E_a}{k_B}\!\left(\frac{1}{T_{ref}}-\frac{1}{T_j}\right)\right]
  \bigl[1+\beta\,\Lambda_s(I,T_w)\bigr]\,
  \bigl(1+\gamma N\bigr)\;}
  \label{eq:master}
\end{equation}

e l'indice di invecchiamento cumulato

\begin{equation}
  A(t) = \int_{0}^{t} \frac{dA}{d\xi}\,d\xi ,
  \qquad
  A \in [0,1],
\end{equation}

con $A = 0$ azionamento nuovo e $A = 1$ fine vita convenzionale. La vita residua
stimata a condizioni operative costanti vale

\begin{equation}
  L_{res} = \frac{1 - A(t)}{\dfrac{dA}{dt}} .
\end{equation}

\section{Implementazione a tempo discreto}

Con periodo di campionamento $T_s$ (tipicamente \SIrange{1}{60}{\second}, poiché
le dinamiche termiche sono lente rispetto al controllo):

\begin{equation}
  A[k+1] = A[k] + T_s\,
  \frac{k_\tau}{L_0}\,
  \left(\frac{I[k]}{I_N}\right)^{m} AF_{Arr}\bigl(T_j[k]\bigr)
  \bigl[1+\beta\Lambda_s\bigr]\bigl(1+\gamma N[k]\bigr).
\end{equation}

\begin{lstlisting}[caption={Pseudo-codice per implementazione embedded.}]
/* eseguito ogni Ts secondi */
Tj_K  = Tj_C + 273.15f;
af_arr = expf( (Ea/kB) * (1.0f/Tref_K - 1.0f/Tj_K) );
af_i   = powf( I/IN, m );
lam    = sigm((Tw - Tw_th)/dTw) * sigm((I - I_th)/dI);
af_th  = 1.0f + beta*lam;
af_n   = 1.0f + gamma*(float)N;
k_tau  = (igbt_type == 1) ? 2.0f : 1.0f;

rate   = (k_tau/L0) * af_i * af_arr * af_th * af_n;   /* 1/h */
A     += rate * (Ts/3600.0f);
if (A > 1.0f) A = 1.0f;
\end{lstlisting}

L'accumulatore $A$ va memorizzato in memoria non volatile con periodicità
compatibile con il numero di cicli di scrittura ammessi.

\section{Parametri suggeriti}

\begin{table}[h]
\centering
\begin{tabular}{lcl}
\toprule
Parametro & Valore & Note \\
\midrule
$E_a$      & \SI{0.5}{\electronvolt} & da qualifica o dati di campo \\
$T_{ref}$  & \SI{100}{\degreeCelsius} & giunzione di riferimento \\
$L_0$      & \SI{1e5}{\hour}          & vita attesa a $T_{ref}$, $I_N$, $\tau=2$, $N=0$ \\
$m$        & 2                        & profilo a carico continuo \\
$T_{w,th}$ & \SI{40}{\degreeCelsius}  & soglia acqua in ingresso \\
$I_{th}$   & $0.8\,I_N$               & soglia di corrente \\
$\beta$    & 1                        & raddoppio del tasso in regime critico \\
$\gamma$   & 0.05                     & penalizzazione per intervento \\
$\Delta T_w$ & \SI{2}{\kelvin}        & regolarizzazione soglia termica \\
$\Delta I$ & $0.05\,I_N$              & regolarizzazione soglia di corrente \\
\bottomrule
\end{tabular}
\caption{Set di parametri di primo tentativo.}
\end{table}

\section{Esempio numerico}

Condizioni: $I = 0.9\,I_N$, $T_w = \SI{45}{\degreeCelsius}$,
$T_j = \SI{125}{\degreeCelsius}$, $N = 2$, modulo di tipo~2.

\begin{align}
  AF_I    &= 0.9^{2} = 0.81 \\
  AF_{Arr}&= \exp\!\left[\frac{0.5}{\num{8.617e-5}}
             \left(\frac{1}{373.15}-\frac{1}{398.15}\right)\right]
             = e^{0.976} = 2.66 \\
  AF_{th} &= 1 + 1\cdot 1 = 2 \quad (T_w > T_{w,th},\ I > I_{th}) \\
  AF_N    &= 1 + 0.05\cdot 2 = 1.10
\end{align}

\begin{equation}
  \frac{dA}{dt} = \frac{1}{\num{1e5}}\cdot 0.81\cdot 2.66\cdot 2\cdot 1.10
  = \SI{4.73e-5}{\per\hour}
  \;\Rightarrow\;
  L \approx \SI{21100}{\hour} \approx \SI{2.4}{yr}.
\end{equation}

Con modulo di tipo~1 il tasso raddoppia e la vita attesa scende a circa
\SI{10600}{\hour}. Rimuovendo la condizione di soglia congiunta (ad esempio
riducendo $T_w$ sotto \SI{40}{\degreeCelsius}) la vita torna a circa
\SI{42200}{\hour}.

\section{Limiti del modello}

\begin{itemize}
  \item Il modello è cumulativo lineare (ipotesi di Miner~\cite{miner1945}):
        non rappresenta interazioni fra meccanismi di degrado né effetti di
        soglia irreversibili.
  \item La legge di Arrhenius sulla sola $T_j$ media non cattura il danno da
        ciclaggio termico rapido; per profili di missione fortemente ciclici va
        affiancato un conteggio rainflow~\cite{downing1982,astm1049} su
        $\Delta T_j$ con curve di Coffin--Manson~\cite{coffin1954,manson1953}.
  \item I coefficienti $k_\tau$, $\beta$, $\gamma$ vanno identificati su dati di
        campo o di prova accelerata; il rapporto $k_1/k_2 = 2$ è qui assunto
        come dato di specifica.
  \item L'accuratezza dipende in modo diretto dalla qualità della stima di
        $T_j$: un errore di \SI{10}{\kelvin} produce, con
        $E_a = \SI{0.5}{\electronvolt}$, un errore sul tasso di circa il
        \SI{30}{\percent}.
\end{itemize}

\begin{thebibliography}{99}

\bibitem{ciappa2002}
M.~Ciappa,
``Selected failure mechanisms of modern power modules,''
\emph{Microelectronics Reliability}, vol.~42, no.~4--5, pp.~653--667, 2002.

\bibitem{lutz2018}
J.~Lutz, H.~Schlangenotto, U.~Scheuermann, R.~De~Doncker,
\emph{Semiconductor Power Devices: Physics, Characteristics, Reliability},
2nd ed., Springer, 2018.

\bibitem{wang2013}
H.~Wang, M.~Liserre, F.~Blaabjerg,
``Toward reliable power electronics: Challenges, design tools, and
opportunities,''
\emph{IEEE Industrial Electronics Magazine}, vol.~7, no.~2, pp.~17--26, 2013.

\bibitem{ma2016}
K.~Ma, H.~Wang, F.~Blaabjerg,
``New approaches to reliability assessment: Using physics-of-failure for
prediction and design in power electronics systems,''
\emph{IEEE Power Electronics Magazine}, vol.~3, no.~4, pp.~28--41, 2016.

\bibitem{held1997}
M.~Held, P.~Jacob, G.~Nicoletti, P.~Scacco, M.-H.~Poech,
``Fast power cycling test of IGBT modules in traction application,''
in \emph{Proc. Int. Conf. Power Electronics and Drive Systems (PEDS)},
Singapore, 1997, pp.~425--430. (modello LESIT)

\bibitem{bayerer2008}
R.~Bayerer, T.~Herrmann, T.~Licht, J.~Lutz, M.~Feller,
``Model for power cycling lifetime of IGBT modules --- various factors
influencing lifetime,''
in \emph{Proc. 5th Int. Conf. Integrated Power Electronics Systems (CIPS)},
Nuremberg, 2008, pp.~1--6.

\bibitem{scheuermann2013}
U.~Scheuermann, R.~Schmidt,
``A new lifetime model for advanced power modules with sintered chips and
optimized Al wire bonds,''
in \emph{Proc. PCIM Europe}, Nuremberg, 2013, pp.~810--817.

\bibitem{avenas2012}
Y.~Avenas, L.~Dupont, Z.~Khatir,
``Temperature measurement of power semiconductor devices by thermo-sensitive
electrical parameters --- A review,''
\emph{IEEE Trans. Power Electronics}, vol.~27, no.~6, pp.~3081--3092, 2012.

\bibitem{miner1945}
M.~A.~Miner,
``Cumulative damage in fatigue,''
\emph{Journal of Applied Mechanics}, vol.~12, no.~3, pp.~A159--A164, 1945.

\bibitem{coffin1954}
L.~F.~Coffin,
``A study of the effects of cyclic thermal stresses on a ductile metal,''
\emph{Transactions of the ASME}, vol.~76, pp.~931--950, 1954.

\bibitem{manson1953}
S.~S.~Manson,
``Behavior of materials under conditions of thermal stress,''
NACA Technical Note~2933, 1953.

\bibitem{downing1982}
S.~D.~Downing, D.~F.~Socie,
``Simple rainflow counting algorithms,''
\emph{International Journal of Fatigue}, vol.~4, no.~1, pp.~31--40, 1982.

\bibitem{astm1049}
ASTM E1049-85,
\emph{Standard Practices for Cycle Counting in Fatigue Analysis},
ASTM International.

\bibitem{jep122}
JEDEC Solid State Technology Association,
\emph{JEP122: Failure Mechanisms and Models for Semiconductor Devices}.

\bibitem{iec60747}
IEC 60747-9,
\emph{Semiconductor devices --- Discrete devices --- Part 9: Insulated-gate
bipolar transistors (IGBTs)}.

\end{thebibliography}

\end{document}
